\subsection{拟伽罗瓦扩张的结构}

命题 13. --- 设 N 是 K 的拟伽罗瓦扩张。我们用 \(N_{r}\) 表示 N 的全部 K-自同构所成的群的不变域，并用 \(N_{s}\) 表示 K 在 N 中的相对可分代数闭包（V，第44页）。则：

\begin{enumerate}
\def\labelenumi{\alph{enumi})}
\item
  \(\mathbf{N}_r\) 是相对 \(p\) -根式闭包（ \(\mathbf{K}\) 在 \(\mathbf{N}\) 中的）（V，第25页）。
\item
  \(N_{s}\) 是 K 的伽罗瓦扩张，且 \(N_{s}\) 的每个 K-自同构以唯一的方式延拓为 N 的一个 \(N_{r}\) -自同构。
\item
  域 \(\mathbf{N}_r\) 和 \(\mathbf{N}_s\) 在 \(\mathbf{K}\) 上线性不相交，并且我们有 \(\mathbf{N} = \mathbf{K}[\mathbf{N}_r \cup \mathbf{N}_s]\) ；换言之， \(\mathbf{N}_r \otimes_{\mathbf{K}} \mathbf{N}_s\) 到 \(\mathbf{N}\) 中的典范同态是同构。
\end{enumerate}

设 \(\Omega\) 为 K 的一个包含 N 作为子扩张的代数闭包（V，第23页，定理2）。 \(\Omega\) 的每个 K-自同构诱导出 N 的一个自同构，因为 N 是拟伽罗瓦的。因此， \(N_r\) 的每个元素在 \(\Omega\) 的 K-自同构群下不变，从而在 K 上是 \(p\) -根式的（V，第53页，推论3）。反之，N 中每个在 K 上 \(p\) -根式的元素显然在 N 的每个 K-自同构下不变，因此属于 \(N_r\) 。这证明了 \(a\) ）。

 \(\Omega\) 的每个 K-自同构将 N 映到 N 中，从而将 \(N_s\) 映到 \(N_s\) 中，所以 \(N_s\) 是 K 的拟伽罗瓦扩张（V，第54页，命题3）。由此可知 \(N_s\) 是 K 的伽罗瓦扩张。 \(N_r \cap N_s\) 的每个元素在 K 上既是可分代数的又是 \(p\) -根式的，因此属于 K（V，第53页，推论3）；因此我们有 \(K_r \cap K_s = K\) 。现在 N 是 \(p\) -根式的（在 \(N_s\) 上）（V，第44页，命题13），且是可分代数的（在 \(N_r\) 上）（V，第56页，定理1），因此既是 \(p\) -根式的又是可分的（在 \(K(N_r \cup N_s)\) 上）。因此我们有 \(N = K(N_r \cup N_s)\) （V，第39页，推论3），且断言 \(b\) ， \(c\) ）由定理5（V，第71页）得出。

推论. --- 设 p 为 K 的特征指数， \(\bar{K}\) 为 K 的一个代数闭包， \(K_{s}\) 为 K 在 \(\bar{K}\) 中的相对可分代数闭包， \(K^{p^{-\infty}}\) 为 K 的完全闭包。则 \(K^{p^{-\infty}} \otimes K_{s}\) 到 \(\bar{K}\) 中的典范同态是同构。

注记. --- 设 R（相应地 S）是 K 的 p-根式（相应地可分代数）扩张。则代数 \(R \otimes_{K} S\) 是一个域：因为 R（相应地 S）同构于 \(K^{p^{-\infty}}\) （相应地 \(K_{s}\) ）的一个子扩张，且只需应用上述推论及 V 第17页的命题1。

